\section{局部正交标架与旋量}

在曲时空里，坐标基通常不正交；而平直时空 Dirac 矩阵的反对易关系写在正交 Minkowski 指标上。旋量也不是普通坐标向量：它按 Lorentz 群的双覆盖群表示变换。因此需要在每点架一组正交基，并说明相邻点的基怎样旋转。

\begin{definition}[正交标架与旋量结构]
在 Lorentz 流形 \((M,g)\) 的一个坐标域内，正交标架 \(e_a=e_a{}^\mu\partial_\mu\) 满足
\(g(e_a,e_b)=\eta_{ab}\)，其对偶余标架 \(e^a=e^a{}_\mu dx^\mu\) 满足
\(g_{\mu\nu}=e^a{}_\mu e^b{}_\nu\eta_{ab}\)。
旋量结构是把正交标架的 \(SO^+(1,3)\) 转换提升到其双覆盖
\(\operatorname{Spin}^+(1,3)\) 的一致选择；下文假定所讨论区域存在这种结构。
\end{definition}

局部标架并非唯一：\(e'_a=\Lambda_a{}^b(x)e_b\) 仍正交。
若 \(\Lambda\) 随 \(x\) 变化，普通偏导 \(\partial_\mu\psi\) 会多出
\((\partial_\mu S)\psi\) 项，不能按旋量表示 \(S(x)\) 协变变换。这正是联络项出现的位置。

\begin{definition}[自旋联络]
由 Levi--Civita 联络定义标架系数
\(\nabla_{\partial_\mu}e_a=\omega_{\mu a}{}^b e_b\)。
度量相容性给 \(\omega_{\mu ab}=-\omega_{\mu ba}\)。
取常数矩阵 \(\gamma^a\) 满足
\(\{\gamma^a,\gamma^b\}=2\eta^{ab}I\)，并令
\(\gamma^{ab}=\tfrac12[\gamma^a,\gamma^b]\)。
旋量协变导数写为
\[
\nabla_\mu^{\mathrm{spin}}\psi
=\partial_\mu\psi-\tfrac14\omega_{\mu ab}\gamma^{ab}\psi,
\]
这里已把 \(\omega_{\mu a}{}^b\) 定义为
\(\nabla_{\partial_\mu}e_a\) 的系数；若改用余标架的系数，
矩阵前的符号也须相应改变。
\end{definition}

这条公式如何工作？令
\(\gamma^\mu(x)=e_a{}^\mu(x)\gamma^a\)，则
\(\{\gamma^\mu,\gamma^\nu\}=2g^{\mu\nu}I\) 由正交标架定义直接得到。
在局部 Lorentz 改变 \(e\mapsto e'\)、\(\psi\mapsto S\psi\) 下，联络矩阵
\(\Omega_\mu=-\tfrac14\omega_{\mu ab}\gamma^{ab}\) 的变换律必须是
\[
\Omega'_\mu=S\Omega_\mu S^{-1}-(\partial_\mu S)S^{-1};
\]
代入可逐项抵消 \(\partial_\mu(S\psi)\) 产生的额外项，得到
\(\nabla'_\mu(S\psi)=S\nabla_\mu\psi\)。
这也是检验联络符号约定是否一致的直接方法。

还可以用 Clifford 关系在一个点上检查系数。因为
\(\gamma^{ab}=\gamma^a\gamma^b-\eta^{ab}I\)，有
\[
[\gamma^{ab},\gamma^c]
=\gamma^a\{\gamma^b,\gamma^c\}
 -\{\gamma^a,\gamma^c\}\gamma^b
=2(\eta^{bc}\gamma^a-\eta^{ac}\gamma^b).
\]
把 \(\Omega_\mu=-\omega_{\mu ab}\gamma^{ab}/4\) 代入，
利用 \(\omega_{\mu ab}=-\omega_{\mu ba}\)，便知
\([\Omega_\mu,\gamma^c]=-\omega_{\mu a}{}^c\gamma^a\)。
另一方面，由
\(\nabla_{\partial_\mu}e_a=\omega_{\mu a}{}^b e_b\)
可写
\(\partial_\mu e_a{}^\nu+\Gamma^\nu_{\mu\rho}e_a{}^\rho
=\omega_{\mu a}{}^b e_b{}^\nu\)。
把两式相加，得到
\[
\partial_\mu\gamma^\nu+\Gamma^\nu_{\mu\rho}\gamma^\rho
+[\Omega_\mu,\gamma^\nu]=0.
\]
这里的加号与上面对旋量导数的减号正好配对；
若只背熟一个联络公式却不检查这一恒等式，就容易把标架约定的
符号混用。

即使空间平直，旋转标架也可有非零联络。在 Euclid 平面的极坐标取
\(e^1=dr,\ e^2=r\,d\varphi\)。Cartan 无挠式
\(de^a+\omega^a{}_b\wedge e^b=0\) 给
\(\omega^2{}_1=d\varphi,\ \omega^1{}_2=-d\varphi\)。
由于 \(d(d\varphi)=0\) 在不含原点的坐标域内成立，曲率仍为零。
所以非零自旋联络系数本身并不说明时空弯曲。

在同一二维例子里，\(\omega_{\varphi12}=1\)、
\(\omega_{\varphi21}=-1\)，所以
\(\Omega_\varphi=-\gamma^{12}/2\)。
取 Euclid Clifford 矩阵满足
\((\gamma^1)^2=(\gamma^2)^2=I\)，则
\(\gamma^2\gamma^{12}=-\gamma^1\)。
因此局部 Dirac 型算子明确写成
\[
\gamma^\mu(\partial_\mu+\Omega_\mu)
=\gamma^1\left(\partial_r+\frac1{2r}\right)
 +\frac{\gamma^2}{r}\partial_\varphi .
\]
\(1/(2r)\) 来自旋转标架，而非平面曲率；
换回 Cartesian 标架后它由 gamma 矩阵的坐标变化承担。

在局部写 Dirac 算子可用
\(\ii\hbar\gamma^\mu(x)\nabla_\mu^{\mathrm{spin}}-mc\)；
具体单位与平直 Dirac 方程的时间、空间坐标约定一致。它同时使用世界指标 \(\mu\) 与局部正交指标 \(a\)，正交标架负责二者之间的转换。

还应说明定义中“双覆盖”一词具体意味着什么。在二维 Euclid 平面，令
\(B=\gamma^1\gamma^2\)，则 Clifford 关系给 \(B^2=-I\)。矩阵
\(S(\theta)=\exp(-\theta B/2)
=\cos(\theta/2)I-B\sin(\theta/2)\)
对向量矩阵 \(v=v_a\gamma^a\) 作共轭
\(v\mapsto S(\theta)vS(\theta)^{-1}\)，得到平面转动；可在
\(\gamma^1\) 上直接求导验证
\(\left.\dd(S\gamma^1S^{-1})/\dd\theta\right|_0
=-\tfrac12[B,\gamma^1]=\gamma^2\)。
但是 \(S(\theta+2\pi)=-S(\theta)\)，而投影到普通向量转动后这两个矩阵给同一个变换。向量绕一周回到自身，旋量在所选提升下多一个负号；绕两周才回到原矩阵。这不是说物理态测量会因整体负号改变，而是说把旋量写成局部标架分量时，不能忽略标架变换的提升选择。

在极坐标例子里，旋转标架沿 \(\varphi\) 绕原点一周，局部提升正是 \(S(\varphi)\)。它在 \(\varphi=0\) 与 \(2\pi\) 两侧相差符号，所以不能在整个穿孔平面上把这一\emph{旋转标架的提升}当作一个单值函数；必须在重叠坐标域用转换函数粘合。穿孔平面本身仍有旋量结构，因为 Cartesian 标架提供了另一种全局描述。区分“这个标架没有全局提升”与“流形不存在旋量结构”，可以避免把局部坐标的符号跳变误判为几何障碍。

\begin{exercise}
对平直 Cartesian 标架求 \(\omega_{\mu ab}\)。再在极坐标例中验证无挠式的 \(a=2\) 分量，并说明为何 \(\omega^2{}_1\ne0\) 不与零曲率矛盾。
\end{exercise}
